\scanpage{6}\mspage{\unsure{\phantom{5}|5}}
% Upper-centre mark: no decipherable numeral first; 5 remains a possible reading of the faint trace.
\[
\underline{\mathrm{Hom}}_{\otimes\mathrm{AUC}}(S^{-1}C,G)\handwritten{\longrightarrow}\underline{\mathrm{Hom}}_{\otimes\mathrm{AUC}}^{S}(C,G)
\]
(où l'exposant $S$ dans le deuxième membre dénote la sous-catégorie pleine formée des $\otimes$-foncteurs AUC $F:C\handwritten{\longrightarrow}G$ tels que $\deleted{P}F(X)$ soit inversib\unsure{\textperiodcentered|-}le pour tout $X\handwritten{\in}\mathrm{Ob}\ S$) est un \uline{isomorphisme} \uline{de catégories}.
% The small mark after inversib stands at the source line end; the following line begins le.
% Its visible remnant is dot-shaped at mid-height; a weak line-break hyphen remains possible.

Nous n'explicitons pas ici la construction, assez naturelle, mais dont l'explicitation pas à pas, \deleted{\rlap{xx}et}\inserted{avec} la vérification de toutes les compatibilités voulues, est fort longue et fastidieuse. Signalons seulement ici que lorsque $S$ est l'ensemble de tous les objets de $C$, $S^{-1}C$ est elle-même une catégorie de Picard (pas nécessairement stricte), qu'on pourra appeler la ``catégorie de Picard enveloppante'' de $C$. C'est sans doute le cas le plus important. Les invariants $\handwritten{\pi}_{0}$ et $\handwritten{\pi}_{1}$ \deleted{\unsure{le|de}}de cette catégorie enveloppante méritent la notation $K^{0}_{\deleted{0}}(C)$ et $K^{\unsure{1|\deleted{0}1|\deleted{1}1}}_{\deleted{1}}(C)$ respectivement. Par exemple, lorsque $C$ est la catégorie des modules projectifs de type fini sur un anneau $A$, \inserted{avec l'opération de somme,} on peut démontrer que ces groupes ne sont autres que les classiques invariants $K^{0}(A)$ et $K^{1}(A)$ étudiés par Grothendieck et par Whitehead-Bass. Le même résultat est sans doute valable pour toute catégorie additive $C$ (munie de \deleted{\rlap{xx}la} l'opération somme), Cela donne une interprétation remarquable des invariants $K^{0}$ et $K^{1}$ en termes d'une ``catégorie stabilisée'' \handwritten{(}de la catégorie additive $C$\handwritten{)}, laquelle est une catégorie de Picard. Signalons seulement que, cette catégorie stabilisée (même dans le cas particulier des modules projectifs de type fini sur l'anneau $A$) n'\deleted{\unsure{est}}\handwritten{ét}ant pratiquement jamais une catégorie de Picard \uline{stricte}, ces considérations semblent rendre peu probable qu'on arrive à construire un objet canonique dans la catégorie dérivée de la catégorie (Ab), qui soit un complexe de chaines satisfaisant $H_{0}(K_{\bullet})=K^{0}$, $H_{1}(K_{\bullet})=K^{1}$ - puisqu'une telle construction (grâce à la théorie de Deligne) nous fournirait précisément une catégorie de Picard \uline{stricte} canonique dont les invariants seraient $K^{0}$ et $K^{1}$.
% At the second K, the primary reading has a live raised 1 with no separately identifiable cancelled upper numeral.
% The preceding line's pi_1 subscript meets this superscript. Cancelled upper 0 or 1 remain alternatives.

On peut difficilement échapper à la question d'une interprétation géométrique analogue des invariants $K^{i}$ supérieurs. Il y a tout lieu de croire qu'une telle interprétation est à\unsure{\textasciidieresis|''} chercher dans la théorie des $n$-catégories de Picard (en un sens convenable) et la construction des $n$-catégories de Picard enveloppantes d'une $\otimes$-catégorie AUC (ou d'une $\otimes$-$n$-catégorie AUC) arbitraire $C$, dont les invariants ``homotopiques'' $\handwritten{\pi}_{i}$ seraient les $K^{i}$ cherchés. De tels développeme\deleted{b}nts obligerai-
% The live text ends here in this scan; no following scan was consulted.

\deleted{Comme \deleted{\unsure{dans}}autre direction de recherche ouverte, suggérée par les résultats esquissés ici, signalons la construction d'une $\otimes$-catégorie de fractions d'une $\otimes$-catégorie AU (sans donnée de commutativité néces-}
% archive sheet number: 5
